\documentclass[11pt]{article}

\usepackage[utf8]{inputenc}
\usepackage[T1]{fontenc}
\usepackage{lmodern}
\usepackage{geometry}
\usepackage{amsmath,amssymb,amsfonts,amsthm,mathtools}
\usepackage{bm}
\usepackage{enumitem}
\usepackage{microtype}
\usepackage{hyperref}
\usepackage{titlesec}
\usepackage{booktabs}
\usepackage{array}
\usepackage{setspace}
\usepackage{mathrsfs}
\usepackage{physics}

\geometry{margin=1in}
\hypersetup{colorlinks=true, linkcolor=black, urlcolor=blue, citecolor=black}
\setlist[enumerate]{topsep=0.4em,itemsep=0.35em}
\setlist[itemize]{topsep=0.3em,itemsep=0.25em}
\titleformat{\section}{\large\bfseries}{\thesection.}{0.5em}{}
\titleformat{\subsection}{\normalsize\bfseries}{\thesubsection.}{0.5em}{}
\setstretch{1.06}

\newcommand{\Aether}{\AE{}ther}
\newcommand{\Aflow}{\AE{}ther-flow}
\newcommand{\TheoryName}{\Aflow{} Interpretation of Relativity}
\newcommand{\TheoryProgram}{\Aether{} / \Aflow{} framework}

\newcommand{\CompletedMetric}{g^{\sharp}}
\newcommand{\MatterSpace}{\Psi}

\newcommand{\Car}{\mathrm{car}}
\newcommand{\Cur}{\mathrm{cur}}
\newcommand{\Hom}{\mathrm{hom}}

\newcommand{\ForSpace}[1]{\mathcal I_{#1}^{\mathrm{for}}}
\newcommand{\PrimShell}{\mathcal S_{\mathrm{prim}}}
\newcommand{\Reduction}[1]{\mathcal R_{#1}}
\newcommand{\CurrentCarrier}[1]{\mathfrak C_{#1}^{\mathrm{cur}}}
\newcommand{\EqCarrier}[1]{\mathfrak C_{#1}^{\mathrm{eq}}}
\newcommand{\MetricOut}{\mathscr G_{\mathrm{out}}}
\newcommand{\MatterOut}{\mathscr T_{\mathrm{out}}}

\newcommand{\SubTarget}[1]{E_{#1}^{\mathrm{sub}}}
\newcommand{\SubPair}[1]{\mathfrak s_{#1}^{\mathrm{sub}}}
\newcommand{\EqnTarget}[1]{Q_{#1}^{\mathrm{eqn}}}
\newcommand{\HiddenSpace}[1]{\mathcal H_{#1}^{\mathrm{hid}}}
\newcommand{\HiddenBody}[1]{D_{#1}^{\mathrm{hid}}}

\newcommand{\ProtoSpace}[1]{\mathcal U_{#1}^{\mathrm{proto}}}
\newcommand{\ProtoBody}[1]{\mathcal B_{#1}^{\mathrm{proto}}}
\newcommand{\ProtoSection}[1]{\Gamma_{#1}^{\mathrm{proto}}}
\newcommand{\ProtoShellSet}[1]{\mathcal Z_{#1}^{\mathrm{proto}}}
\newcommand{\ProtoZeroCore}[1]{\mathcal Z_{#1}^{0,\mathrm{proto}}}
\newcommand{\ProtoSplit}[1]{\Phi_{#1}^{\mathrm{proto}}}
\newcommand{\ProtoCharge}[1]{\mathcal Q_{#1}^{\mathrm{proto}}}
\newcommand{\ProtoEmbed}[1]{\zeta_{#1}^{\mathrm{proto}}}
\newcommand{\ProtoKernelCh}[1]{\widetilde K_{#1}^{0,\mathrm{proto}}}
\newcommand{\ProtoStateComp}[1]{P_{#1}^{\mathrm{proto,co}}}

\newcommand{\PreProtoNucleus}{\mathfrak Q^{\mathrm{nuc}}}
\newcommand{\PreProtoSpace}[1]{\mathcal V_{#1}^{\mathrm{preproto}}}
\newcommand{\PreProtoBody}[1]{\mathcal B_{#1}^{\mathrm{preproto}}}
\newcommand{\PreProtoProj}[1]{\Pi_{#1}^{\mathrm{preproto}}}
\newcommand{\PreProtoProtoMin}[1]{\mathfrak f_{#1}^{\mathrm{preproto,proto}}}
\newcommand{\PreProtoSection}[1]{\Gamma_{#1}^{\mathrm{preproto}}}
\newcommand{\PreProtoFunctional}[1]{\mathscr W_{#1}^{\mathrm{preproto}}}
\newcommand{\PreProtoShellSet}[1]{\mathcal Z_{#1}^{\mathrm{preproto}}}

\newcommand{\PreNucCore}{\mathfrak R^{\mathrm{core}}}
\newcommand{\PreNucSpace}[1]{\mathcal W_{#1}^{\mathrm{prenuc}}}
\newcommand{\PreNucBody}[1]{\mathcal B_{#1}^{\mathrm{prenuc}}}
\newcommand{\PreNucProj}[1]{\Pi_{#1}^{\mathrm{prenuc}}}
\newcommand{\PreNucNucleusMin}[1]{\mathfrak f_{#1}^{\mathrm{prenuc,nuc}}}
\newcommand{\PreNucShellSet}[1]{\mathcal Z_{#1}^{\mathrm{prenuc}}}

\newcommand{\PreCoreTransportCore}{\mathfrak S^{\mathrm{tr}}}
\newcommand{\PreCoreTrBody}[1]{\mathcal B_{#1}^{\mathrm{precore,tr}}}
\newcommand{\PreCoreTrProj}[1]{\widehat{\Pi}_{#1}^{\mathrm{precore}}}
\newcommand{\PreCoreTrLift}[1]{\mathfrak f_{#1}^{\mathrm{precore,tr}}}
\newcommand{\PreCoreTrShell}[1]{\mathcal Z_{#1}^{\mathrm{precore,tr}}}

\newcommand{\SubPreCorePackage}{\mathfrak U^{\mathrm{subpre}}}
\newcommand{\SubPreCoreSpace}[1]{\mathcal L_{#1}^{\mathrm{subpre}}}
\newcommand{\SubPreCoreBody}[1]{\mathcal B_{#1}^{\mathrm{subpre}}}
\newcommand{\SubPreCoreProj}[1]{\Pi_{#1}^{\mathrm{subpre}}}
\newcommand{\SubPreCoreTransportMin}[1]{\mathfrak f_{#1}^{\mathrm{subpre,tr}}}
\newcommand{\SubPreCoreShellSet}[1]{\mathcal Z_{#1}^{\mathrm{subpre}}}

\newcommand{\SubPreImageCore}{\mathfrak U^{\mathrm{img}}}
\newcommand{\SubPreImgBody}[1]{\mathcal B_{#1}^{\mathrm{subpre,img}}}
\newcommand{\SubPreImgProj}[1]{\widehat{\Pi}_{#1}^{\mathrm{subpre}}}
\newcommand{\SubPreImgShell}[1]{\mathcal Z_{#1}^{\mathrm{subpre,img}}}

\newtheorem{definition}{Definition}
\newtheorem{proposition}{Proposition}
\newtheorem{remark}{Remark}
\newtheorem{corollary}{Corollary}
\newtheorem{theorem}{Theorem}

\title{\textbf{The \TheoryName{}}\\[0.4em]
\large Primitive-Reservoir Observer-Reduction\\
\large Exact-Shell-Transport-Core-Generating Substrate Sub-Pre-Core Dynamics\\
\large Collapse to an Observer-Relevant Sub-Pre-Core Exact-Shell Transport-Image Core}
\author{}
\date{}

\begin{document}

\maketitle

\begin{abstract}
The current primitive-reservoir observer-reduction frontier already moved the live burden below the observer-relevant pre-core exact-shell transport core itself: the active sub-pre-core theorem derived that transport core from a deeper exact-shell-transport-core-generating substrate sub-pre-core package. Under the recorded routing safeguard and the benchmark one-metric upstream compression rule, however, the honest next move at that level is not another same-output deeper-origin relay that leaves the sub-pre-core package unchanged. The sharper benchmark-facing gain is to prove a genuinely smaller theorem-level collapse strictly inside that sub-pre-core package itself.

The present note takes that route. It isolates the \emph{observer-relevant sub-pre-core exact-shell transport-image core}: the minimizer-image compact body over the recorded pre-core exact-shell transport body, the restricted projection back to that body, the exact-shell transport image of the recorded pre-core transport shell, and benchmark faithfulness. The full exact sub-pre-core shell is then shown to be forced by that transport image, and ambient sub-pre-core state-space directions outside the minimizer image do not add independent benchmark-facing freedom once the downstream transport-core target is held fixed.

The benchmark boundary remains fixed throughout. The one operative metric is still exactly \(\CompletedMetric\), the ordinary matter route is unchanged, there is no second operative metric, the low-energy matter law is not enlarged, and no stronger promotion language is licensed. The positive result is therefore narrower than another deeper relay and sharper than the full sub-pre-core package: the live burden now collapses to a smaller observer-relevant sub-pre-core exact-shell transport-image core. What remains open is to derive or justify that smaller core from still more primitive substrate structure, or else prove a still sharper collapse strictly inside it.
\end{abstract}

\section{Introduction}

The active derivational program of \TheoryName{} is again constrained by two routing facts. First, the current active-primary theorem already removed the observer-relevant pre-core exact-shell transport core as the primitive stopping point by deriving it from a more primitive exact-shell-transport-core-generating substrate sub-pre-core package. Second, the repository has already fixed a safeguard against recursive depth drift: once a benchmark-relevant core is in hand, another deeper-origin relay that preserves that same live burden is no longer honest front-line progress merely by being deeper.

That safeguard matters here. The next honest move below the current sub-pre-core frontier is therefore not to spend a manuscript on another same-output relay unless that relay genuinely changes the benchmark-facing burden. The sharper default at current scope is instead to ask whether the full recorded sub-pre-core package contains ambient freedom that is benchmark-neutral and can be collapsed away without reopening the benchmark exact-closure package.

The present note answers that question affirmatively. Inside the recorded exact-shell-transport-core-generating substrate sub-pre-core package, the ambient state-space presentation and off-image body directions are not an independent observer-relevant burden. Once attention is restricted to the minimizer image over the recorded pre-core exact-shell transport body and the exact-shell transport image of the recorded transport shell, the omitted ambient directions are forced or benchmark-neutral.

\paragraph{Framework and claim boundary.}
\TheoryProgram{} denotes the broader \Aether{} / \Aflow{} framework built on the claims that the \Aether{} is the underlying four-dimensional substrate of reality, the \Aflow{} is its intrinsic ordered motion, observed three-dimensional space is the local experiential slice of that deeper substrate, S-time is the experienced order of change arising from matter, light, and the \Aflow{}, and observed expansion is the three-dimensional appearance of deeper four-dimensional ordered motion. Within that framework, \TheoryName{} denotes the exact-closure theory in which the effective gravitational dynamics are adopted to be exactly Einsteinian with universal matter coupling. In the active sequence, \emph{adoption} means use of that established relativistic dynamics without claiming substrate derivation, while \emph{derivation} is reserved for a first-principles recovery from explicit substrate variables.

The present note stays strictly inside that boundary. It does not introduce a second operative metric, it does not enlarge the low-energy matter law, and it does not reopen older screened layers as if they were again the live burden. Its positive role is narrower and benchmark facing: compress the live sub-pre-core burden to the part of the package that still matters observer-relevantly.

\section{Fixed Inputs}

\begin{definition}[Recorded primitive continuation data]
\label{def:fixed}
Fix the following continuation data from the active primitive-reservoir observer-reduction line.
\begin{enumerate}
    \item For each sector label \(\sigma\in\{\Car,\Cur,\Hom\}\), a primitive forcing shell
    \[
    \ForSpace{\sigma},
    \]
    a visible reduction map
    \[
    \Reduction{\sigma}:\ForSpace{\sigma}\to X_{\sigma},
    \]
    and shell-level current and equilibrium carriers
    \[
    \CurrentCarrier{\sigma}:\ForSpace{\sigma}\to Z_{\sigma}^{\mathrm{cur}},
    \qquad
    \EqCarrier{\sigma}:\ForSpace{\sigma}\to Z_{\sigma}^{\mathrm{eq}}.
    \]
    \item The product shell
    \[
    \PrimShell
    \equiv
    \ForSpace{\Car}\times \ForSpace{\Cur}\times \ForSpace{\Hom}.
    \]
    \item The recorded metric and ordinary-matter outputs
    \[
    \MetricOut:\PrimShell\to\{\text{observer metrics}\},
    \qquad
    \MatterOut:\PrimShell\times\MatterSpace\to\{\text{observer stress tensors}\},
    \]
    satisfying
    \begin{equation}
    \MetricOut(\mathbf w)=\CompletedMetric,
    \qquad
    \MatterOut(\mathbf w,\psi)
    =
    T_{\mu\nu}^{\mathrm{matter}}[\CompletedMetric,\psi]
    \label{eq:fixedoutputs}
    \end{equation}
    for every \(\mathbf w\in\PrimShell\) and every matter configuration \(\psi\in\MatterSpace\).
\end{enumerate}
\end{definition}

\begin{definition}[Recorded proto-germ exact-shell target]
\label{def:proto}
Fix \(\sigma\in\{\Car,\Cur,\Hom\}\). The active line already fixes:
\begin{enumerate}
    \item the same substrate target and pairing
    \[
    \SubTarget{\sigma},
    \qquad
    \SubPair{\sigma}:
    \SubTarget{\sigma}\times\SubTarget{\sigma}\to\mathbb R;
    \]
    \item a hidden fiber space \(\HiddenSpace{\sigma}\) with compact hidden body \(\HiddenBody{\sigma}\subset\HiddenSpace{\sigma}\);
    \item a finite-dimensional proto-germ state space \(\ProtoSpace{\sigma}\), a compact convex proto-germ body \(\ProtoBody{\sigma}\subset\ProtoSpace{\sigma}\), a smooth proto-germ restriction section
    \[
    \ProtoSection{\sigma}:\ProtoSpace{\sigma}\to\SubTarget{\sigma},
    \]
    and the exact proto-germ shell
    \[
    \ProtoShellSet{\sigma}
    \equiv
    \bigl(\ProtoSection{\sigma}\bigr)^{-1}(0)\cap\ProtoBody{\sigma};
    \]
    \item a closed embedded proto zero core \(\ProtoZeroCore{\sigma}\subset\ProtoShellSet{\sigma}\) and a split diffeomorphism
    \[
    \ProtoSplit{\sigma}:
    \ProtoZeroCore{\sigma}\times\HiddenBody{\sigma}
    \xrightarrow{\cong}
    \ProtoShellSet{\sigma};
    \]
    \item a hidden charge
    \[
    \ProtoCharge{\sigma}:\HiddenSpace{\sigma}\to\EqnTarget{\sigma},
    \]
    a smooth observer-compatible exact proto-germ zero-response realization
    \[
    \ProtoEmbed{\sigma}:\ForSpace{\sigma}\hookrightarrow\ProtoZeroCore{\sigma},
    \]
    and a fixed-data solvability / coercive package on \(\ProtoZeroCore{\sigma}\), recorded through \(\ProtoKernelCh{\sigma}\) and \(\ProtoStateComp{\sigma}\).
\end{enumerate}
\end{definition}

\begin{definition}[Recorded current transport nucleus]
\label{def:nucleus}
Fix \(\sigma\in\{\Car,\Cur,\Hom\}\). The recorded current transport nucleus \(\PreProtoNucleus\) for sector \(\sigma\) consists of:
\begin{enumerate}
    \item the pre-proto-germ state space \(\PreProtoSpace{\sigma}\), compact convex body \(\PreProtoBody{\sigma}\subset\PreProtoSpace{\sigma}\), projection
    \[
    \PreProtoProj{\sigma}:\PreProtoSpace{\sigma}\to\ProtoSpace{\sigma},
    \]
    and proto section
    \[
    \PreProtoProtoMin{\sigma}:\ProtoBody{\sigma}\to\PreProtoBody{\sigma}
    \]
    satisfying
    \[
    \PreProtoProj{\sigma}\bigl(\PreProtoBody{\sigma}\bigr)=\ProtoBody{\sigma},
    \qquad
    \PreProtoProj{\sigma}\circ\PreProtoProtoMin{\sigma}
    =
    \mathrm{id}_{\ProtoBody{\sigma}};
    \]
    \item the restriction section
    \[
    \PreProtoSection{\sigma}:\PreProtoSpace{\sigma}\to\SubTarget{\sigma},
    \]
    the induced functional
    \[
    \PreProtoFunctional{\sigma}(q)
    \equiv
    \SubPair{\sigma}\bigl(\PreProtoSection{\sigma}(q),\PreProtoSection{\sigma}(q)\bigr),
    \]
    and the exact shell
    \[
    \PreProtoShellSet{\sigma}
    \equiv
    \bigl(\PreProtoSection{\sigma}\bigr)^{-1}(0)\cap\PreProtoBody{\sigma},
    \]
    with
    \[
    \PreProtoSection{\sigma}\circ\PreProtoProtoMin{\sigma}
    =
    \ProtoSection{\sigma},
    \qquad
    \PreProtoProj{\sigma}\big|_{\PreProtoShellSet{\sigma}}
    :
    \PreProtoShellSet{\sigma}
    \xrightarrow{\cong}
    \ProtoShellSet{\sigma};
    \]
    \item benchmark faithfulness, meaning preservation of the benchmark outputs \eqref{eq:fixedoutputs}, no second operative metric, no enlarged low-energy matter law, and no premature promotion from adoption to completed derivation.
\end{enumerate}
\end{definition}

\begin{definition}[Recorded observer-relevant pre-nucleus exact-shell core]
\label{def:core}
Fix \(\sigma\in\{\Car,\Cur,\Hom\}\). The observer-relevant pre-nucleus exact-shell core \(\PreNucCore\) for sector \(\sigma\) consists of:
\begin{enumerate}
    \item the pre-nucleus state space \(\PreNucSpace{\sigma}\), compact convex body \(\PreNucBody{\sigma}\subset\PreNucSpace{\sigma}\), affine surjection
    \[
    \PreNucProj{\sigma}:\PreNucSpace{\sigma}\to\PreProtoSpace{\sigma},
    \]
    and fiber minimizer
    \[
    \PreNucNucleusMin{\sigma}:\PreProtoBody{\sigma}\to\PreNucBody{\sigma}
    \]
    satisfying
    \[
    \PreNucProj{\sigma}\bigl(\PreNucBody{\sigma}\bigr)=\PreProtoBody{\sigma},
    \qquad
    \PreNucProj{\sigma}\circ\PreNucNucleusMin{\sigma}
    =
    \mathrm{id}_{\PreProtoBody{\sigma}};
    \]
    \item a closed embedded exact pre-nucleus shell
    \[
    \PreNucShellSet{\sigma}\subset\PreNucBody{\sigma}
    \]
    satisfying
    \[
    \PreNucNucleusMin{\sigma}\bigl(\PreProtoShellSet{\sigma}\bigr)
    \subset
    \PreNucShellSet{\sigma},
    \qquad
    \PreNucProj{\sigma}\big|_{\PreNucShellSet{\sigma}}
    :
    \PreNucShellSet{\sigma}
    \xrightarrow{\cong}
    \PreProtoShellSet{\sigma};
    \]
    \item benchmark faithfulness in the sense of Definition \ref{def:nucleus}.
\end{enumerate}
\end{definition}

\begin{definition}[Recorded pre-core exact-shell transport core]
\label{def:transportcore}
Fix \(\sigma\in\{\Car,\Cur,\Hom\}\). The recorded observer-relevant pre-core exact-shell transport core \(\PreCoreTransportCore\) for sector \(\sigma\) consists of:
\begin{enumerate}
    \item a compact transport-core body
    \[
    \PreCoreTrBody{\sigma},
    \]
    an affine projection
    \[
    \PreCoreTrProj{\sigma}:\PreCoreTrBody{\sigma}\to\PreNucBody{\sigma},
    \]
    and a smooth transport-core lift
    \[
    \PreCoreTrLift{\sigma}:\PreNucBody{\sigma}\to\PreCoreTrBody{\sigma}
    \]
    satisfying
    \[
    \PreCoreTrProj{\sigma}\bigl(\PreCoreTrBody{\sigma}\bigr)=\PreNucBody{\sigma},
    \qquad
    \PreCoreTrProj{\sigma}\circ\PreCoreTrLift{\sigma}
    =
    \mathrm{id}_{\PreNucBody{\sigma}};
    \]
    \item a closed embedded exact transport shell
    \[
    \PreCoreTrShell{\sigma}\subset\PreCoreTrBody{\sigma}
    \]
    satisfying
    \[
    \PreCoreTrLift{\sigma}\bigl(\PreNucShellSet{\sigma}\bigr)
    \subset
    \PreCoreTrShell{\sigma},
    \qquad
    \PreCoreTrProj{\sigma}\big|_{\PreCoreTrShell{\sigma}}
    :
    \PreCoreTrShell{\sigma}
    \xrightarrow{\cong}
    \PreNucShellSet{\sigma};
    \]
    \item benchmark faithfulness in the sense of Definition \ref{def:core}.
\end{enumerate}
\end{definition}

\begin{remark}
Definitions \ref{def:fixed}--\ref{def:transportcore} are fixed inputs. The one operative metric, the ordinary matter route, the fixed proto-germ exact-shell target, the current transport nucleus, the recorded pre-nucleus exact-shell core, and the recorded pre-core exact-shell transport core are not reopened here. The present note asks only which part of the sub-pre-core package still carries independent benchmark-facing burden once those downstream data are already fixed.
\end{remark}

\section{Full Sub-Pre-Core Package}

\begin{definition}[Exact-shell-transport-core-generating substrate sub-pre-core dynamics]
\label{def:subpre}
Fix \(\sigma\in\{\Car,\Cur,\Hom\}\). An exact-shell-transport-core-generating substrate sub-pre-core dynamics package \(\SubPreCorePackage\) for sector \(\sigma\) consists of the following data.
\begin{enumerate}
    \item A finite-dimensional Euclidean sub-pre-core state space \(\SubPreCoreSpace{\sigma}\), a compact convex sub-pre-core body
    \[
    \SubPreCoreBody{\sigma}\subset\SubPreCoreSpace{\sigma}
    \]
    with nonempty interior, an affine surjection
    \[
    \SubPreCoreProj{\sigma}:\SubPreCoreSpace{\sigma}\to\PreCoreTrBody{\sigma},
    \]
    and a smooth sub-pre-core minimizer
    \[
    \SubPreCoreTransportMin{\sigma}:\PreCoreTrBody{\sigma}\to\SubPreCoreBody{\sigma}
    \]
    satisfying
    \[
    \SubPreCoreProj{\sigma}\bigl(\SubPreCoreBody{\sigma}\bigr)=\PreCoreTrBody{\sigma},
    \qquad
    \SubPreCoreProj{\sigma}\circ\SubPreCoreTransportMin{\sigma}
    =
    \mathrm{id}_{\PreCoreTrBody{\sigma}}.
    \]
    \item A closed embedded exact sub-pre-core shell
    \[
    \SubPreCoreShellSet{\sigma}\subset\SubPreCoreBody{\sigma}
    \]
    satisfying
    \[
    \SubPreCoreTransportMin{\sigma}\bigl(\PreCoreTrShell{\sigma}\bigr)
    \subset
    \SubPreCoreShellSet{\sigma},
    \qquad
    \SubPreCoreProj{\sigma}\big|_{\SubPreCoreShellSet{\sigma}}
    :
    \SubPreCoreShellSet{\sigma}
    \xrightarrow{\cong}
    \PreCoreTrShell{\sigma}.
    \]
    \item Benchmark faithfulness, meaning preservation of the benchmark outputs \eqref{eq:fixedoutputs}, no second operative metric, no enlarged low-energy matter law, and no premature promotion from adoption to completed derivation.
\end{enumerate}
\end{definition}

\begin{remark}
Definition \ref{def:subpre} is the full current sub-pre-core package. The present note does not spend its move on another deeper relay beneath it. It asks instead whether the full package is already minimal once the recorded transport-core target is fixed.
\end{remark}

\section{Observer-Relevant Sub-Pre-Core Exact-Shell Transport-Image Core}

\begin{definition}[Observer-relevant sub-pre-core exact-shell transport-image core]
\label{def:imagecore}
Fix \(\sigma\in\{\Car,\Cur,\Hom\}\) and a benchmark-faithful sub-pre-core package in the sense of Definition \ref{def:subpre}. The \emph{observer-relevant sub-pre-core exact-shell transport-image core} \(\SubPreImageCore\) for sector \(\sigma\) consists of:
\begin{enumerate}
    \item the compact transport-image body
    \[
    \SubPreImgBody{\sigma}
    \equiv
    \SubPreCoreTransportMin{\sigma}\bigl(\PreCoreTrBody{\sigma}\bigr)
    \subset
    \SubPreCoreBody{\sigma},
    \]
    together with the restricted projection
    \[
    \SubPreImgProj{\sigma}
    \equiv
    \SubPreCoreProj{\sigma}\big|_{\SubPreImgBody{\sigma}}
    :
    \SubPreImgBody{\sigma}\to\PreCoreTrBody{\sigma};
    \]
    \item the exact-shell transport image
    \[
    \SubPreImgShell{\sigma}
    \equiv
    \SubPreCoreTransportMin{\sigma}\bigl(\PreCoreTrShell{\sigma}\bigr)
    \subset
    \SubPreCoreShellSet{\sigma},
    \]
    with the restricted shell projection
    \[
    \SubPreImgProj{\sigma}\big|_{\SubPreImgShell{\sigma}}
    :
    \SubPreImgShell{\sigma}\to\PreCoreTrShell{\sigma};
    \]
    \item benchmark faithfulness in the sense of Definition \ref{def:subpre}.
\end{enumerate}
\end{definition}

\begin{remark}
Definition \ref{def:imagecore} is intentionally smaller than Definition \ref{def:subpre}. It keeps only the minimizer image over the recorded pre-core exact-shell transport body together with the exact-shell transport it already determines. The ambient sub-pre-core state space and off-image body directions are not taken as independent core data.
\end{remark}

\section{Collapse of the Full Sub-Pre-Core Package}

\begin{proposition}[The transport-image body carries the body-level content]
\label{prop:body}
For every benchmark-faithful sub-pre-core package, \(\SubPreImgBody{\sigma}\) is compact and the restricted projection
\[
\SubPreImgProj{\sigma}:\SubPreImgBody{\sigma}\to\PreCoreTrBody{\sigma}
\]
is a homeomorphism with inverse \(\SubPreCoreTransportMin{\sigma}\).
\end{proposition}

\begin{proof}
The body \(\PreCoreTrBody{\sigma}\) is compact by Definition \ref{def:transportcore}, and \(\SubPreCoreTransportMin{\sigma}\) is continuous by Definition \ref{def:subpre}(1). Therefore
\[
\SubPreImgBody{\sigma}
=
\SubPreCoreTransportMin{\sigma}\bigl(\PreCoreTrBody{\sigma}\bigr)
\]
is compact.

By Definition \ref{def:subpre}(1),
\[
\SubPreCoreProj{\sigma}\circ\SubPreCoreTransportMin{\sigma}
=
\mathrm{id}_{\PreCoreTrBody{\sigma}}.
\]
Hence
\[
\SubPreImgProj{\sigma}\circ\SubPreCoreTransportMin{\sigma}
=
\mathrm{id}_{\PreCoreTrBody{\sigma}}.
\]
Conversely, every \(y\in\SubPreImgBody{\sigma}\) has the form
\[
y=\SubPreCoreTransportMin{\sigma}(w)
\qquad
\text{for some }
w\in\PreCoreTrBody{\sigma},
\]
so
\begin{align*}
\SubPreCoreTransportMin{\sigma}\bigl(\SubPreImgProj{\sigma}(y)\bigr)
&=
\SubPreCoreTransportMin{\sigma}\bigl(\SubPreCoreProj{\sigma}(y)\bigr) \\
&=
\SubPreCoreTransportMin{\sigma}\bigl(\SubPreCoreProj{\sigma}(\SubPreCoreTransportMin{\sigma}(w))\bigr) \\
&=
\SubPreCoreTransportMin{\sigma}(w) \\
&=
y.
\end{align*}
Thus \(\SubPreImgProj{\sigma}\) and \(\SubPreCoreTransportMin{\sigma}\) are mutually inverse continuous maps between compact Hausdorff spaces, hence homeomorphisms.
\end{proof}

\begin{proposition}[The full exact sub-pre-core shell is forced by the transport image]
\label{prop:shell}
For every benchmark-faithful sub-pre-core package,
\[
\SubPreImgShell{\sigma}
=
\SubPreCoreShellSet{\sigma},
\]
and the restricted shell projection
\[
\SubPreImgProj{\sigma}\big|_{\SubPreImgShell{\sigma}}
:
\SubPreImgShell{\sigma}\xrightarrow{\cong}\PreCoreTrShell{\sigma}
\]
is a diffeomorphism with inverse \(\SubPreCoreTransportMin{\sigma}\big|_{\PreCoreTrShell{\sigma}}\).
\end{proposition}

\begin{proof}
Definition \ref{def:subpre}(2) already gives the inclusion
\[
\SubPreCoreTransportMin{\sigma}\bigl(\PreCoreTrShell{\sigma}\bigr)
\subset
\SubPreCoreShellSet{\sigma},
\]
which is the inclusion
\[
\SubPreImgShell{\sigma}\subset\SubPreCoreShellSet{\sigma}.
\]

For the reverse inclusion, let \(z\in\SubPreCoreShellSet{\sigma}\). Since
\[
\SubPreCoreProj{\sigma}\big|_{\SubPreCoreShellSet{\sigma}}
:
\SubPreCoreShellSet{\sigma}\xrightarrow{\cong}\PreCoreTrShell{\sigma}
\]
is a diffeomorphism by Definition \ref{def:subpre}(2), the point
\[
w\equiv\SubPreCoreProj{\sigma}(z)
\]
lies in \(\PreCoreTrShell{\sigma}\). Again by Definition \ref{def:subpre}(2),
\[
\SubPreCoreTransportMin{\sigma}(w)\in\SubPreCoreShellSet{\sigma}.
\]
But
\[
\SubPreCoreProj{\sigma}\bigl(\SubPreCoreTransportMin{\sigma}(w)\bigr)=w=\SubPreCoreProj{\sigma}(z),
\]
and the shell projection is injective because it is a diffeomorphism. Therefore
\[
z=\SubPreCoreTransportMin{\sigma}(w)\in\SubPreImgShell{\sigma}.
\]
Hence
\[
\SubPreCoreShellSet{\sigma}\subset\SubPreImgShell{\sigma},
\]
so equality holds. The final statement follows because the full shell projection is already a diffeomorphism and its inverse on \(\PreCoreTrShell{\sigma}\) is exactly \(\SubPreCoreTransportMin{\sigma}\big|_{\PreCoreTrShell{\sigma}}\).
\end{proof}

\begin{proposition}[Downstream exact-shell transport factors through the transport-image core]
\label{prop:downstream}
For every benchmark-faithful sub-pre-core package, the recovery of the recorded pre-core exact-shell transport core and all downstream shell transport to the fixed pre-nucleus exact-shell core, transport nucleus, and proto-germ exact-shell target factor through \(\SubPreImageCore\). More precisely,
\[
\PreCoreTrProj{\sigma}\circ\SubPreImgProj{\sigma}
:
\SubPreImgBody{\sigma}\to\PreNucBody{\sigma}
\]
is a surjection with section
\[
\SubPreCoreTransportMin{\sigma}\circ\PreCoreTrLift{\sigma}
:
\PreNucBody{\sigma}\to\SubPreImgBody{\sigma},
\]
and the composite shell transport
\[
(\PreCoreTrProj{\sigma}\circ\SubPreImgProj{\sigma})\big|_{\SubPreImgShell{\sigma}}
:
\SubPreImgShell{\sigma}
\xrightarrow{\cong}
\PreNucShellSet{\sigma}
\]
is a diffeomorphism. Consequently
\[
(\PreNucProj{\sigma}\circ\PreCoreTrProj{\sigma}\circ\SubPreImgProj{\sigma})\big|_{\SubPreImgShell{\sigma}}
:
\SubPreImgShell{\sigma}
\xrightarrow{\cong}
\PreProtoShellSet{\sigma}
\]
and
\[
(\PreProtoProj{\sigma}\circ\PreNucProj{\sigma}\circ\PreCoreTrProj{\sigma}\circ\SubPreImgProj{\sigma})\big|_{\SubPreImgShell{\sigma}}
:
\SubPreImgShell{\sigma}
\xrightarrow{\cong}
\ProtoShellSet{\sigma}
\]
are also diffeomorphisms.
\end{proposition}

\begin{proof}
By Proposition \ref{prop:body},
\[
\SubPreImgProj{\sigma}\circ\SubPreCoreTransportMin{\sigma}
=
\mathrm{id}_{\PreCoreTrBody{\sigma}}.
\]
Composing with the recorded transport-core data of Definition \ref{def:transportcore}(1) yields
\[
(\PreCoreTrProj{\sigma}\circ\SubPreImgProj{\sigma})\bigl(\SubPreImgBody{\sigma}\bigr)
=
\PreCoreTrProj{\sigma}\bigl(\PreCoreTrBody{\sigma}\bigr)
=
\PreNucBody{\sigma}
\]
and
\[
(\PreCoreTrProj{\sigma}\circ\SubPreImgProj{\sigma})\circ
\SubPreCoreTransportMin{\sigma}\circ\PreCoreTrLift{\sigma}
=
\PreCoreTrProj{\sigma}\circ\PreCoreTrLift{\sigma}
=
\mathrm{id}_{\PreNucBody{\sigma}}.
\]

For shell transport, Proposition \ref{prop:shell} gives a diffeomorphism
\[
\SubPreImgProj{\sigma}\big|_{\SubPreImgShell{\sigma}}
:
\SubPreImgShell{\sigma}\xrightarrow{\cong}\PreCoreTrShell{\sigma},
\]
while Definition \ref{def:transportcore}(2) gives a diffeomorphism
\[
\PreCoreTrProj{\sigma}\big|_{\PreCoreTrShell{\sigma}}
:
\PreCoreTrShell{\sigma}\xrightarrow{\cong}\PreNucShellSet{\sigma}.
\]
Their composition is therefore a diffeomorphism from \(\SubPreImgShell{\sigma}\) to \(\PreNucShellSet{\sigma}\). Composing once more with the shell diffeomorphism of Definition \ref{def:core}(2) yields the displayed diffeomorphism to \(\PreProtoShellSet{\sigma}\), and composing finally with the shell diffeomorphism of Definition \ref{def:nucleus}(2) yields the displayed diffeomorphism to \(\ProtoShellSet{\sigma}\).
\end{proof}

\begin{proposition}[Ambient sub-pre-core directions are benchmark neutral at current scope]
\label{prop:neutral}
For every benchmark-faithful sub-pre-core package, the ambient sub-pre-core state space \(\SubPreCoreSpace{\sigma}\) and the off-image body directions
\[
\SubPreCoreBody{\sigma}\setminus\SubPreImgBody{\sigma}
\]
do not add independent benchmark-facing freedom beyond \(\SubPreImageCore\).
\end{proposition}

\begin{proof}
By Proposition \ref{prop:body}, the full body-level transport relevant to the recorded pre-core exact-shell transport core is already fixed by the transport-image body \(\SubPreImgBody{\sigma}\), the restricted projection \(\SubPreImgProj{\sigma}\), and its inverse \(\SubPreCoreTransportMin{\sigma}\). By Proposition \ref{prop:shell}, the full exact sub-pre-core shell is not additional data beyond the transport image of the recorded transport shell; it is exactly that image. By Proposition \ref{prop:downstream}, all downstream transport to the recorded pre-nucleus exact-shell core, the current transport nucleus, and the fixed proto-germ exact-shell target factors through this smaller image core. Finally, benchmark faithfulness is already part of Definition \ref{def:imagecore}(3) and preserves the fixed outputs \eqref{eq:fixedoutputs}. Therefore nothing benchmark-facing at the present frontier requires the ambient sub-pre-core state-space presentation or body points outside the minimizer image as additional independent data.
\end{proof}

\begin{proposition}[The benchmark package remains fixed]
\label{prop:benchmark}
Every observer-relevant sub-pre-core exact-shell transport-image core preserves the same benchmark one-metric package \eqref{eq:fixedoutputs}. In particular, the operative metric remains exactly \(\CompletedMetric\), the ordinary matter route is unchanged, there is no second operative metric, and the low-energy matter law is not enlarged.
\end{proposition}

\begin{proof}
This is part of benchmark faithfulness in Definition \ref{def:imagecore}(3). The present note only removes benchmark-neutral ambient sub-pre-core freedom from the live burden. It does not alter the benchmark exact-closure package.
\end{proof}

\section{Main Theorem}

\begin{theorem}[Exact-shell-transport-core-generating substrate sub-pre-core dynamics collapse to an observer-relevant sub-pre-core exact-shell transport-image core]
\label{thm:main}
Fix the primitive continuation data of Definition \ref{def:fixed}, the recorded proto-germ exact-shell target of Definition \ref{def:proto}, the recorded current transport nucleus of Definition \ref{def:nucleus}, the recorded observer-relevant pre-nucleus exact-shell core of Definition \ref{def:core}, the recorded observer-relevant pre-core exact-shell transport core of Definition \ref{def:transportcore}, and a benchmark-faithful sub-pre-core package in the sense of Definition \ref{def:subpre}. Then, for each sector \(\sigma\in\{\Car,\Cur,\Hom\}\), the live benchmark-facing burden inside the current sub-pre-core frontier collapses from the full sub-pre-core package \(\SubPreCorePackage\) to the smaller observer-relevant sub-pre-core exact-shell transport-image core \(\SubPreImageCore\). More precisely:
\begin{enumerate}
    \item the body-level content of the full sub-pre-core package is already carried by the transport-image body \(\SubPreImgBody{\sigma}\) and the restricted projection \(\SubPreImgProj{\sigma}\) by Proposition \ref{prop:body};
    \item the full exact sub-pre-core shell is forced to equal the transport image of the recorded pre-core transport shell by Proposition \ref{prop:shell};
    \item all downstream recovery of the recorded pre-core exact-shell transport core, the pre-nucleus exact-shell core, the current transport nucleus, and the fixed proto-germ exact-shell target factors through \(\SubPreImageCore\) by Proposition \ref{prop:downstream};
    \item ambient sub-pre-core directions outside that image core are benchmark neutral at current scope by Proposition \ref{prop:neutral};
    \item the benchmark boundary remains fixed by Proposition \ref{prop:benchmark};
    \item consequently the remaining honest burden below the previous sub-pre-core frontier is no longer the full exact-shell-transport-core-generating substrate sub-pre-core dynamics package \(\SubPreCorePackage\) as such. What remains is to derive or justify the sharper image core \(\SubPreImageCore\) from still more primitive substrate structure, or else prove a still smaller theorem-level collapse strictly inside that image core.
\end{enumerate}
\end{theorem}

\begin{proof}
Item (1) is Proposition \ref{prop:body}. Item (2) is Proposition \ref{prop:shell}. Item (3) is Proposition \ref{prop:downstream}. Item (4) is Proposition \ref{prop:neutral}. Item (5) is Proposition \ref{prop:benchmark}. Item (6) is the routing consequence of the previous items together with the active safeguard against counting another same-output deeper-origin relay as new front-line progress when the live benchmark-relevant datum is unchanged.
\end{proof}

\begin{corollary}[What must be done next]
\label{cor:next}
After Theorem \ref{thm:main}, the next honest benchmark-facing manuscript must either:
\begin{enumerate}
    \item derive or justify the observer-relevant sub-pre-core exact-shell transport-image core \(\SubPreImageCore\) from still more primitive substrate structure in a way that changes benchmark-facing status;
    \item identify a genuinely smaller theorem-level observer-relevant datum strictly inside \(\SubPreImageCore\) and prove a sharper collapse to it;
    \item do either of the previous two while preserving the same benchmark one-metric package and the same claim boundary.
\end{enumerate}
Another same-output deeper-origin relay that preserves this image core, another re-expansion back to the full sub-pre-core package as if its screened ambient directions were still live, another reopening of the already-screened larger pre-core package as if it were again the active burden, or any move that introduces a second operative metric, enlarges the ordinary matter law, or uses stronger promotion language does not satisfy the present burden. If a quantum continuation is pursued, it matters benchmark-facingly only insofar as it derives this same image core, collapses it further, or closes a benchmark derivation gate in a genuinely new way.
\end{corollary}

\begin{proof}
Theorem \ref{thm:main} shows that the full sub-pre-core package is no longer the minimal benchmark-facing datum at this stage. Therefore a subsequent manuscript must improve on the smaller image core rather than merely restating the larger package.
\end{proof}

\section{Conclusion}

The positive result of this note is a disciplined compression step inside the current sub-pre-core frontier. The previous active-primary theorem already removed the observer-relevant pre-core exact-shell transport core as the primitive stopping point by deriving it from a deeper sub-pre-core package. The present theorem then takes the remaining honest internal route available there: it shows that the full sub-pre-core package is itself not the minimal benchmark-facing datum.

What survives as the live burden is the observer-relevant sub-pre-core exact-shell transport-image core. That sharper datum consists of the minimizer-image compact body over the recorded pre-core exact-shell transport body, the restricted projection back to that body, the exact-shell transport image of the recorded pre-core transport shell, and benchmark faithfulness. The full exact sub-pre-core shell is forced to equal that transport image, and ambient sub-pre-core directions outside the minimizer image do not add independent benchmark-facing freedom once the downstream transport-core target is fixed.

This is the honest next burden after the present theorem. A new primary manuscript must derive or justify that sharper image core from still more primitive substrate structure, or else collapse it further, while preserving the same benchmark one-metric package and the same claim boundary. The current result does not yet complete a first-principles derivation of GR from substrate microphysics, but it does narrow the live derivational burden one level further on the active line.

\end{document}
